% English draft based on sections/04_response_channels_ja.tex
\section{Response Channels and Silent Worlds}

\S{}3 showed that $L_F$ has no internal stopping condition. This chapter asks how that non-closure appears in different operating environments. The distinction between responsive and silent worlds is not ontological. It is a threshold phenomenon.

\subsection{Responsive and Silent Regimes}

A responsive world is an environment in which an inferential output receives a return. The agent asks, acts, hypothesizes, or experiments, and the world responds. In C-geometry, the gradient forms a closed loop: the output changes the world, and the changed world returns as feedback.

A silent world is an environment in which such return does not close. The agent can output inference, but no corresponding response comes back. Ultimate reasons, fixed past events, and unresolved future contingencies often appear in this mode.

Formally, a responsive regime has significant mutual constraint in both directions, $H_{ij}$ and $H_{ji}$. A silent regime has a broken or negligible return channel, approximately $H_{ji}\approx 0$.

\subsection{Silence Is Not Defect}

Silence is not a defect in the world. It is not merely ignorance waiting to be solved. Some regions do not return the kind of response required to close the inferential loop. If responsive-world strategies are applied to silent-world regions without threshold mobility, inference can enter infinite regress.

\subsection{Threshold Phenomenon}

Responsive and silent regimes are phases determined by thresholds. The same environment can appear responsive or silent depending on the position of the boundary set by parameters such as inferential velocity, resolution, and stopping hardness. The threshold mobility $\mu_\theta$ is the ability to move this boundary.

\begin{figure}[H]
\centering
\includegraphics[width=\linewidth]{fig2_responsive_silent_threshold.pdf}
\caption{Responsive and silent regimes. These are not two ontological kinds of world, but threshold-dependent phases determined by the boundary moved by $\mu_\theta$.}
\label{fig:part2-responsive-silent}
\figurenote{Both panels share the same environment and the same C-landscape. What changes is the threshold boundary that determines whether a response loop closes.}{In the responsive phase, external mutual constraint supplies local stopping. In the silent phase, this supply is cut off, creating pressure to externalize stopping elsewhere.}
\end{figure}

\subsection{Infinite Regress in Silent Worlds}

If an agent is fixed in a silent regime, three conditions hold: the environment does not return a stopping response, justification can only be advanced by inference, and stopping cannot be derived from inference itself. Under these conditions, inference lacks a stopping resource and continues indefinitely.

Infinite regress and branching explosion are two aspects of the same structural fact. Internal non-fixedness means the trajectory does not become still; branching means the trajectory is not single. In silent regimes, the two combine. The result is a doubly unstable dynamics.

\subsection{The Preview of $L_R$}

Non-stopping inference leaves traces. The trajectory records where it passed, which branches it attempted, and where it diverged. These traces form logs. Recursive reference to such logs is the beginning of the log layer $L_R$.

This is the bridge to \S{}5. The failure of $L_F$ to stop is not only a dead end. It is also productive: because inference does not stop, traces accumulate; because traces accumulate, recursive reference becomes possible; because recursive reference becomes possible, $L_E$ can emerge under external mutual constraint.

\bigskip
